\chapter{What has the battlefield become?}

UxS have made movement one of the most dangerous tools on the battlefield. In the \textit{grey zone}, exposure to observation dictates risk. This chapter examines how democratised surveillance and armed UxS have reshaped warfare.

%\section{Peculiarities of the Russo--Ukrainian War (559 words)}
While context-dependent, the Russo--Ukrainian War exposes lessons which are applicable to militaries worldwide. \textcite[9--13]{MOLLOY_2024} and \textcite[62--63]{ZABRODSKYI_2022} describe the adoption of drones since the outbreak of hostilities. The war's positional \parencite{watling_2024_panel}, attritional \parencite{JONES_2023,WATLING_2023A,WATLING_2023B,WATLING_2025B,WATLING_2025n,ZABRODSKYI_2022} character is the product of mutually-denied airspace operating under the shadow of nuclear deterrence, rather than manoeuvre's obsolescence \parencite[5:20--10:20]{BRONK_2022a} \parencite[25:00--27:00]{WATLING_2022C} \parencite[22--23,~31]{hvizda_2025}. 

However, Western militaries must resist uncritical lesson importation. \textcite[1--2]{Watling_book} argues that abstract demonstrations of capability often obscure the enabling requirements necessary to make them plausible in practice. Some analysts consider that drones have rendered manned airpower irrelevant \parencite{BARNO_2024}, while others caution that armoured and air warfare remain crucial \parencite{barrie_2021,BRONK_2025,SHAIKH_2020}. \textcite{stu_lyle_participant_11} assessed that combined-arms manoeuvre could exploit the positional dependence of drone teams, dislocating their positions and reducing their battlefield effect. \textcite[5:00--6:00]{CORNISH_2024} situates drone warfare within a broader taxonomy of ideational, kinetic, technological and cognitive warfare, cautioning against overemphasising the technological dimension --- echoing earlier scholars' perspectives. The offence--defence balance currently favours the defender\footnote{Though both sides continue to adapt; see \textcite{JONES_2023} for battlefield evolution.}.

%Others push back against manoeuvrist orthodoxy, contending that positional and attritional fighting --- far from an aberration --- represents a recurring operational precondition for manoeuvre, one Western doctrine has been too quick to dismiss \parencite[12:58--15:15]{FOX_2021,FOX_2025,Halem_2024,Simoens_24,KING_2023}. Following \textcite{MOLLOY_2024}'s comprehensive evidence-based research, she contends that ``uncrewed systems are disrupting [... and ...] have a place in modern warfare''. \textcite[p62--63]{ZABRODSKYI_2022} elucidates the character of the positional warfare in Ukraine --- one dominated by UAS.

\posscite{laboaratory} contention of Ukraine as a ``laboratory of new tech'' warrants careful examination. Militaries are often accused of preparing to fight the last war, yet there is strategic value in seeking to close peacetime gaps in RCP. States cannot know in advance whether war will merely expose peacetime weaknesses or punish them. The safest peacetime strategy is to seek parity --- underscoring the relevance of the ongoing war.


\textcite[15]{watling_2025_10} considers pervasive ISR to be ``the single most transformative element of warfare in Ukraine''. \textcite{bob_tollast_participant_12} cautions against dismissing UxS as a fleeting phenomenon and against overcorrecting through radical force redesign, noting that comparatively low cost equipment and TTPs may suffice. Meanwhile \textcite{stu_lyle_participant_11} cautions against bureaucratic inertia, citing \posscite{paine_1776_common_sense} ``a long habit of not thinking a thing wrong gives it a superficial appearance of being right''.


%\section{Environmental Constraints}
Weather and terrain create windows for action which light infantry must learn to exploit. \gls{rc} UAS\index{Uncrewed Aerial Systems} usually operate within the microwave band\index{Microwave Band} which is attenuated by liquid water \parencite{vollmer}. Hence wet environments, rain and humidity reduce their range. Greater battlespace transparency heightens the importance of concealment. While poor weather has long enabled concealed movement, modern tactics place a premium on conditions inhibiting small rotary-wing UAS \parencite[3:15--4:56]{civdiv_6}. Fog obscures visual detection, while headwinds degrade endurance \parencite[0:10--0:56]{civdiv_3} --- exacerbated by their circa $25$-minute lithium-polymer battery life \parencite{civdiv_1,dan_participant_1}. Furthermore, water vapour degrades uncooled thermal imaging. Consequently, UAS employed around the grey zone remain highly weather-dependent \parencite[10]{kunertova_2024_2024_08,miller_2026}. Preferential times for action include poor weather, dusk and dawn --- reminiscent of ISIS' \textit{thunderstorm assaults}.

%\subection{Terrain\index{Terrain} (307 words)}\label{terrain}
Broken terrain and clutter --- which is separated on the order of the signal's wavelength --- limits line of sight\index{Line Of Sight}, reduces the range of RC drones \parencite[31,53]{neron_bancel_2024_05_27} and is compounded further by the opacity inherent in the land environment. This can be mitigated somewhat by employing UAS to rebroadcast signals. Since RC range is highly sensitive to terrain, \gls{fpvowa}\index{First Person View One-Way Attack Drone!Terminal phase altitude at long range} remain vulnerable even with relay support: after reaching the target area, they may be unable to descend to attack altitude without dropping the control link. \gls{term:casd} with higher operating altitudes and less stringent line-of-sight requirements, are somewhat less exposed to this constraint. Hence, close terrain reduces the persistence and pervasiveness of RC UAS\index{Radio Control Uncrewed Aerial Systems} \parencite[1:00--1:57]{civdiv_4}. Furthermore, for similar reasons, dense terrain reduces the range of EW jamming\index{Electronic Warfare|seealso{CUAS}}\index{Jamming|see{Electronic Warfare}}. Unlike the flat nature of Ukrainian steppe where EW transmitter range is optimised, close terrain will necessitate careful planning and judicious selection of target areas of interest\index{Target Areas Of Interest}. 

While ISR\index{Intelligence Surveillance And Reconnaissance} is increasingly effective on the modern battlefield, the trend exemplified in Ukraine can be misleading if generalised \parencite{bob_tollast_participant_12}. The flat steppe and limited vegetation are exceptionally transparent to ISR collection, maximising detection range and hindering safe movement. Closer and hillier terrain provide opportunities to remain undetected by channeling surveillance to named areas of interest\index{Named Areas Of Interest}. For other areas, concealment\index{Concealment} may once again offer meaningful mitigation against both observation\index{Observation} and engagement\index{Engagement}, reducing the near-equivalence\index{Detection}\index{Engagement} seen in Ukraine.




%\subection{Limits of drone observation (181 words)}

Assertions that modern technologies have eliminated the \textit{fog of war}\index{Fog of War} and produced a transparent battlespace are not new. They have consistently proven overly optimistic \parencite[49]{ALACH_2008}. \textcite{DAVIS_e_2025} cautions against concluding that battlefield transparency will ever be fully realised. \textcite{bob_tollast_participant_12} argues that Ukraine’s open terrain offers near-ideal conditions for ISR drones, making it misleading to attribute transparency solely to drones rather than to the terrain they exploit.

Human factors\index{Human Factors} also constrain drone observation. Pilots\index{Drone Pilots} perceive the battlespace\index{Battlespace} through a two-dimensional\index{Two-Dimensional} visual interface\index{Visual Interface}, lacking depth perception\index{Depth Perception} and peripheral awareness\index{Peripheral Awareness}. This limits their ability to interpret texture\index{Texture} and subtle environmental changes\index{Environmental Change} --- emphasising movement\index{Movement} as the dominant detection cue. The \gls{slls}\index{SLLS} drill provides contextual understanding\index{Senses!Contextual Understanding} which remote sensing cannot replicate \parencite[4:35--5:32]{civdiv_6}.



%\section{Persistent ISR in the Contested Grey Zone (1022 words)}
\index{Democratisation!Of transparency}As battlefield observation becomes cheaper and more widespread, close-air lethality increases and exposure becomes more costly \parencite[15]{hvizda_2025} and \parencite{neron_bancel_2024_05_27}. \textcite{SBS_Jan_2023_infinopen} depicts how light forces without CUAS are vulnerable to attrition by ISR-corrected fires before they can close with the enemy. \textcite{kraken_2022_12_30} demonstrates that Ukraine's early-war tactics followed well-understood Western combined-arms doctrine. These methods are no longer survivable. However, Western forces have shown marginal and uneven adaptation. The French Foreign Legion’s \textit{Gévaudan 26} exercise is illustrative: despite reported FPVD integration \parencite{ministere_des_armees_2026_gevaudan_26_legionnaires_combat}, the observed use of tight formations, lack of CUAS sentries, exposed movement and force concentration remains inconsistent with drone warfare \parencite{tf1_info_2026_legion_etrangere_entrainement}. NATO’s Exercise Hedgehog further demonstrates the disproportionate lethality which skilled UAS operators can generate \parencite{melchior_2026}. 

\textcite{darragh_participant_4} noted that in 2022, infantry could still ``scurry into a wood line'' and wait for surveillance drones to leave; they ``didn't have persistence''. By 2023--24, as \textcite{fred_participant_9} describes, the conflict had shifted: both sides were employing layered \gls{term:isrd}, while loitering munitions and FPV-OWADs converted detection into almost immediate lethality \parencite{bob_tollast_participant_12,JONES_2023,WATLING_2023A,WATLING_2024A,ZABRODSKYI_2022}. Ukrainian tactics have evolved under the experience of frontline conditions often involving up to 25 UAS tracking a single formation \parencite[32]{WATLING_2023B}. Sections can expect four ISRDs to be simultaneously overhead --- often with friendly drones indistinguishable from the enemy's --- creating fratricidal and other operational frictions \parencite{fred_participant_9}. Both belligerents field ISRDs at platoon level, with commensurate EW\index{Electronic Warfare}\index{Counter UAS!Electronic Warfare} countermeasures on each side \parencite[14]{hvizda_2025}. 



This logic collectively embodies what this thesis terms \textit{movement as timed exposure}\index{Movement as Timed Exposure} --- routes must be pre-planned\index{Route Planning} and remotely reconnoitred; soldiers move quickly in dispersed packets\index{Dispersion}\index{Movement!Dispersion} between cover and concealment; pauses are required to check for drones; and delay increases the likelihood of detection and engagement.

 


Reports such as \textcite[22]{MOLLOY_2024} and veterans such as \textcite{dan_participant_1} describe the {grey zone or \textit{zone of contestation}\index{Zone of Contestation|seealso{Grey Zone}}\index{Grey Zone|seealso{Zone of Contestation}}, where mutual sensor-fires parity makes concentration prohibitively dangerous. This view enjoys broad scholarly support: \textcite[23]{neron_bancel_2024_05_27}, \textcite{tollast_2025_11}, \textcite[17:21--22:37]{WATLING_2023_book_launch}, \textcite[~53--58]{Watling_book}, \textcite[~2--3]{watling_2025_10} and \textcite[~53--55]{ZABRODSKYI_2022}. Extending $15$--$30$ kilometres past the FLOT --- or \textit{zero line} --- \parencite[5]{Tyshchenko_2026_03_20_ukraine_bunker_story,watling_2025_10}, the grey zone has supplanted the traditional \textit{close--deep--rear} model with zones of \textit{opportunity--contestation--risk} \parencite[~53--58]{Watling_book}. Shaping operations must now position objectives within the friendly zone of opportunity prior to their assault\index{Zone of Opportunity}\index{Zone of Contestation}\index{Zone of Risk}.


Cheap UAS have delivered persistence and rapidly scalable strike at the small-unit level \parencite[4:11--5:12]{civdiv_3}, \parencite[23]{kunertova_2024_2024_08}, \parencite[~26--28]{neron_bancel_2024_05_27} and \parencite[2]{watling_2025_10}. \textcite{tollast_2025_11} emphasises that ISRDs at lower echelons are central to transparency. Their utility in light infantry assaults is striking: they enable real-time adjusting of fires, battle damage assessment and decision making.

The mass arrival of Russian FPV-OWADs\index{First Person View One-Way Attack Drone!Effect on Russo--Ukrainian War} ``completely turn[ed] the tide'' \parencite{darragh_participant_4,fred_participant_9}. Detection--to--engagement times have fallen to as little as three seconds \parencite{stu_lyle_participant_11}. ISRDs also compress the kill-chain, providing target coordinates directly to fire controllers \parencite[29]{MOLLOY_2024,neron_bancel_2024_05_27}. Movement corridors and logistics routes become ``predictable and targetable'', degrading tempo and increasing attrition \parencite{WATLING_2024A}.

 




On the ground, this persistence is experiential: ``the buzz of rotors overhead is constant'', with fibre-optic cables and netted roads physically marking a battlespace \parencite{miller_2026}. For example, the \textcite{225_oshp_2025_12_24_FPV_HUNT_INf} depicts the first-person view of a motorised infantry mortar team being effectively targeted by FPVDs. Veterans now describe the Russo--Ukrainian battlespace as: ``impossible not to be seen'' \textcite{dan_participant_1}, ``you will get seen, you will get attacked'' \parencite{garand_thumb_CUAS_2026_3_15}.

Infantry have adapted through dispersed, ``packeted'' infiltration, deliberate SLLS\index{SLLS}\index{Drone Halt} \textit{drone halts} and strict signature discipline \parencite{dan_participant_1,dan_participant_6,fred_participant_9,ZABRODSKYI_2022} and \parencite[8:11--8:23]{civdiv_3}. Command posts and drone teams likewise minimise electromagnetic emissions and disperse to survive fires \parencite{Watling_book}. Such dispersion, however, fundamentally disaggregates the force: small elements operating under signature discipline cannot be centrally controlled. Indeed, \textcite[9]{WATLING_2025B} report section defences with as much as $200$ m frontage\index{Mission Command!Defensive frontage}\index{Mission Command!Dispersion}. This places a premium on subordinate initiative and pre-delegated authority through mission command.


Greater battlefield transparency imposes a structural trade-off between survivability and the ability to generate effects \parencite[4:36--6:06]{WATLING_2023_book_launch}. Survivability can be enhanced through concealment or hardened positions, but these constrain effective action. Subterranean positions illustrate this: their entrances must be secured against drones and infantry, constraining manoeuvre \parencite{civdiv_2}. \textcite{darragh_participant_4} describes assaulting positions where entrances had been welded shut --- prioritising protection at the cost of tactical flexibility --- and apparently armed with an ``unlimited amount of grenades''.

%\section{Armed UAS in the Contested Grey Zone (486 words)}
Armed UAS have produced an asymmetric cost imposition which fundamentally changes the economics of attrition for light forces. Lethality in the Russo--Ukrainian War has exhibited a marked switch from artillery to drones \parencite{zafra_2024_03_26_drone_combat_ukraine}, where a purported $80$--$90$ percent of casualties are caused by UAS \parencite{hartog_2026_01_26_drones_casualties_ukraine,TCAG_Brad,zelenskyy_2026_03_17} --- which is consistent with the expectations of this research for light forces. This reflects a substantial asymmetry in cost and effect. \textcite{magyar_telegram_MAR_22} reports that \textit{Magyar's Birds}\index{Magyar's Birds} and associated UAS units account for over a third of verified Russian losses while comprising just two percent of the army's headcount --- at a reported cost of \$878 per enemy soldier killed --- a striking asymmetric cost. Their comparatively low cost reduces the penalty for inefficiency, enabling redundant expenditure of drones on a single target. $10''$ CASDs and $5''$ FPV-OWADs can be expected to deliver light-mortar sized warheads at ranges of seven kilometres ($3.5$ return) for circa \$500 \parencite[2:08--3:20]{civdiv_8}\index{First Person View Drone!Cost}. While radio-controlled systems remain the most common \parencite[0:04--0:18]{civdiv_5}, fibre-optic tethering provides resilience against jamming and extends operational range to tens of kilometres. \textcite[10]{MOLLOY_2024} identifies the prevalence of small drones as the most frequently cited theme among participants. As these systems operate predominantly at the tactical level alongside infantry, this reinforces the need to integrate uncrewed systems into the planning and preparation of infantry forces. 

%\subection{FPV-OWADs}\index{First Person View One-Way Attack Drone!Targeting light infantry}

While \textcite[15]{watling_2025_10} acknowledges the ``utility'' of FPV-OWADs, two of his suggested limitations merit closer scrutiny. First, the ease of protecting positions against them appears overstated for light infantry: positions are small (vulnerable to blast) and require firing ports, creating unavoidable vulnerabilities which FPVDs can exploit. Second, their purported vulnerability to interception underestimates the effects of their small size, high speed and frequent nocturnal employment (constituting a particular light infantry vulnerability). \posscite{Leppanen_2025} field demonstration reinforces the central claim of this thesis: that light infantry positions are fundamentally exposed to FPV-OWADs under realistic conditions. In a fashion analogous to human ambushes, \gls{waiters} attempt to optimise concealment with observation of a target area --- possibly for weeks \parencite{Stepanenko_2025b}\index{First Person View One-Way Attack Drone!``Waiters''}. While they are vulnerable to being pre-seen and destroyed --- particularly where the drone has a single camera --- detection is not usually a warning but rather an immediate prelude to strike.

%\subection{Dropper Drones}
Dropper drones --- typified by systems such as AFUA's \textit{Baba Yaga} --- represent the most functionally versatile UxS in the contemporary battlespace. Unlike FPV-OWADs, they are reusable, payload-versatile and multi-role: delivering CAS, acting as motherships, emplacing mines along logistics corridors and conducting resupply \parencite[8]{watling_2025_10,azov_media_20251111_escape_impossible_drones_mines,garand_thumb_CUAS_2026_3_15}. \textcite[12]{watling_2025_10} assesses CASDs as among the deadliest tools in Ukraine's arsenal, with innovations including double-drops and magazine-fed release systems increasingly emulating crewed aircraft \parencite{murray_tuga,taylor_ukr}.
Their significance extends beyond lethality. By assuming the most exposure-intensive tasks --- last-mile resupply and route interdiction --- dropper drones displace human movement along predictable, targetable routes. This directly mitigates \textit{movement as timed exposure}. While scholars emphasise the enduring primacy of indirect fires on grounds of responsiveness, payload and range, CASDs contest each criterion: they deliver repeated precision strikes with substantial explosive payloads and operate from positions already near the FLOT --- compressing effective range without the signature or positional constraints of tube systems.

%\section{UGSs (1238 words)}
UGS do not replace infantry; they displace the most dangerous human movement, thereby sustaining combat power under persistent observation. The Russo--Ukrainian War has accelerated the adaptation of UGS for logistics, reconnaissance and fire support --- particularly since the failure of Ukraine's 2023 offensive \parencite{Kofman_Lee_2023_09_04} and \parencite[4--5]{watling_2025_10}. Dedicated units are emerging which employ UGS for assault, mine-laying, resupply and CASEVAC, offering light infantry a means to sustain combat effectiveness while mitigating exposure to increasingly lethal and transparent battlefields \parencite{bendett_2024_06_17,KIRICHENKO_2022A,third_army_corps_2025_10_23,watling_2025_10}. Where personnel are constrained, drones allow fewer soldiers to be exposed to risk while extracting greater utility from limited human resources over time. \posscite[ch.~9]{Watling_book} conclusions regarding the future offensive and logistical use of UGS by light infantry have largely been realised.

%\subection{Piloting and Navigation}

UGS differ fundamentally from UAS since they must negotiate the physical complexity of the terrain rather than overfly it. Constrained by undulations, debris, gradient and soil conditions, tipping or immobilisation is a recurring vulnerability. Although their low profile may reduce detection and enhance survivability \parencite{david_participant_2}, ground navigation imposes difficulties absent from aerial systems. This is consistent with the ``hidden difficulty'' described by \textcite{Rivero_2026_0_03}, reported by others \parencite{manley_2024_04_06_ground_drone_trench_attack} and experienced by the author.

The implication is that effective UGS employment is more terrain and weather dependent than aerial systems, demands a level of operator judgement and local adaptation which cannot be scripted from above which reinforces the premium on decentralised control examined in Chapter 2.

%\subection{Logistics}

UGS are emerging as a sustainment technology for conditions in which human movement itself has become the principal source of vulnerability. As \textcite[33]{MOLLOY_2024} identifies, drones displace rather than replace humans. It is this displacement which sustains combat power under persistent observation --- a logic now operationalised at scale in Ukraine --- as illustrated by \textcite{mod_russia_20260315_UGS}. There UGS assume the final, most dangerous leg of sustainment after crews deliver them by vehicle to a forward release point, beyond which human drivers and stretcher-bearers can no longer operate safely \parencite{david_participant_2}.

At Pokrovsk --- a heavily contested Russo--Ukrainian War axis --- the displacement was already extensive. \textcite{Shatalov_2025} (a frontline commander), \textcite{third_army_corps_2025_10_23}, \parencite{david_participant_2} and \textcite{Rivero_2026_0_03} report that approximately $80$--$90$ percent of forward resupply and casualty evacuation are conducted by UGS. This reflects what \textcite[57]{Watling_book} identifies as the core utility of robotics for dispersed light infantry --- sustaining them without requiring the movement that would expose them to detection and targeting. It should be noted that CASDs are frequently re-rolled as logistics drones \parencite[18]{watling_2025_10}.


The operational stakes were illustrated near Kostiantynivka\index{Kostiantynivka UGS Resupply}\index{UGS!Resupply at Kostiantynivka}. UGS delivered medical supplies, a stretcher and warm clothing under simultaneous artillery and FPV attack, with a second vehicle completing the evacuation of wounded soldiers after repeated human attempts had failed or generated further casualties \parencite{Kirichenko_2026B,bendett_2024_06_17}. UGS therefore do more than reduce risk. In a battlespace where grey zone approaches are saturated with ISRDs \parencite{david_participant_2}, they make ``last mile resupply'' and casualty recovery possible where manned movement has become operationally infeasible \parencite[40:00--41:52]{lyle_23_interview}. Indeed, \textcite[18:50--28:00]{taylor_ukr} reports UGS range greater than $20$ kilometres. 

The displacement of human movement by machines represents a structural shift in how combat power is sustained. By transferring risk from personnel to platforms, robotics mitigate the vulnerability of light infantry to bypass, positional warfare and encirclement. This logic extends beyond high-intensity conflict, offering a means to reduce the \textit{targetable action} inherent in resupplying forward operating bases and isolated positions in inter-positional operations such as crisis management and counter-insurgency.


%\subection{Offence}

Interview-based reporting from a Ukrainian UGS company commander indicates that armed ground drones are used for support-by-fire, harassing, breaching defences, sabotage, ambushes and reinforcing infantry assaults at the tactical level \parencite{devdroid_2026_drones_change_warfare_interview,suprun_2025_09_05_nc13_nrk_interview,third_army_corps_2025_2_26}. Ukrainian payloads from tens to hundreds of kilograms are reported possible \parencite{United24Media_2024_10_25_UGS}. UGS can deliver substantial explosive payloads for breaching defences and destroying known positions, illustrated by \textcite{Fedorov_2024_4_2,drone_surrender_1}.

The complexity and difficulty of piloting and navigation likely limit the utility of UGS during fast-paced manoeuvre. However, their value in breaching, suppression and attrition is evident, echoing earlier uses of uncrewed systems by ISIS \parencite{tveteras_2022_islamic_state_vehicle_industry} and bearing comparison to the static logic of First World War mine warfare. Their potential as disaggregated support-by-fire platforms is likewise significant\index{UAS!Support-by-fire tasks}\index{UAS!Defensive fire support}, as illustrated by \textcite[fig.~4]{watling_2025_10}. This suggests that the tripod-mounted machine-gun as the primary light infantry direct fire weapon may face structural redundancy in this role. The skilful use of ISRD-corrected machine guns in the indirect-fire role has also been observed. UGS are suited to urban operations, where they can conduct reconnaissance by fire and provide heavy direct-fires, while mitigating human exposure in these challenging roles \parencite[151]{stu_lyle_participant_11,Watling_book}. They are also useful for reconnaissance during building clearance, although extensive rubble may favour UAS. Future quadrupedal robots may offer a replacement for pack animals in difficult terrain, enhancing light infantry sustainment while increasing combat effectiveness and lethality. This reflects a structural shift in which firepower can be projected without exposing personnel, with operators controlling systems from protected positions or armoured vehicles.

In July 2025, Ukraine’s 3rd Separate Assault Brigade\index{3rd Separate Assault Brigade (Ukraine)} reportedly induced the surrender of Russian troops using FPV-OWADs\index{First Person View One-Way Attack Drone!Surrender of soldiers} and UGS, suggesting that UxS may influence not only physical survivability but also an adversary’s will to fight \parencite{drone_surrender_1,drone_surrender_2,drone_surrender}. The psychological effects of human--machine combat remain underexplored.


%\subection{Defence}
Armed UGS provide a means to project defensive firepower while reducing human exposure, functioning as  platforms for infantry support \parencite{david_participant_2}. Ukrainian experience demonstrates their growing institutionalisation, with specialised units such as the NC13\index{NC13 UGS Strike Company}\footnote{A strike UGS company of the 3rd Assault Brigade.} employing robotic systems for support, ambush and defensive tasks. It reflects a broader effort to replace personnel in the most dangerous sectors of the front \parencite{devdroid_2026_drones_change_warfare_interview,suprun_2025_09_05_nc13_nrk_interview}. 

In practice, armed UGS have demonstrated the ability to assume limited defensive roles, including replacing infantry positions and maintaining continuous fire over key terrain for extended periods, thereby relieving pressure on manpower-constrained units \parencite{Militarnyi_Shumlianskyi_2025_11_15}. However, their utility remains constrained by vulnerability to UAS, limited endurance, communication fragility and the requirement for human support to sustain operations, maintain UGS and ultimately secure terrain \parencite{david_participant_2}. Consequently, while they enhance defensive capacity and redistribute risk, they do not replace infantry and may be better suited to covering gaps along quieter sectors of the front --- providing nuance to recent rhetorical claims \textcite{loh_rboto,tretii_armiiskyi_korpus_2025_12_22_ground_drone_holds_position}. Moreover, their relative cost compared to aerial drones imposes compromise: while they reduce risk to personnel, they are not readily expendable and must be employed selectively within a broader human--machine defensive system. While CASEVAC UGS and dispersed infantry seek to survive by reducing their signature and becoming less attractive targets, logistics and defensive UGS may instead constitute \textit{more interesting} high-payoff targets\index{High Payoff Target} for the adversary \parencite{david_participant_2} and \parencite[40:00--42:00]{lyle_23_interview}.












%\section{Avoiding Drones}
Infantry have adapted by turning movement into a calculated exercise in timed exposure, using dispersion, signature discipline and deliberate pauses to survive. Practitioner evidence suggests that drone avoidance requires a fundamental shift regarding how light infantry manage exposure, time and signature. As \textcite{garand_thumb_CUAS_2026_3_15} emphasises, traditional camouflage principles remain central. UAS detection is primarily movement-based; identification via shape, shine, shadow, spacing or silhouette is far less reliable without AI \parencite[5:00--6:00]{civdiv_1}. However, camouflage\index{Camouflage!As personal survivability equipment} is best considered a component of one's survivability equipment \parencite{Hall_2021_camouflage_platform}, comparable to body armour. The democratisation\index{Democratisation!Of precision strike} of precision strike described by \textcite{stu_lyle_participant_11} and \textcite[14]{hvizda_2025} results in camouflage shaping whether a soldier is detected --- which is the first stage in the adversary's kill-chain \parencite{Hemez_2017}. \textcite[3]{watling_2025_10} observes Russian forces employing thermal sheeting --- adaptation not merely to visual detection but to multi-spectral surveillance, a shift further illustrated by \textcite{ukrainian_navy_2025_12_7_camo} and \textcite{United24Media_camouflage_2026_2_28}.

Crucially, the tempo of the drone-enabled kill-chain forces a shift away from static concealment towards dynamic survivability. This often requires minimising direct exposure to the sky, though doing so may create tensions with dispersion when concealment is limited. While earlier Russian practice --- taken from blogger Filatov --- emphasised freezing to avoid detection, emerging evidence suggests this is increasingly untenable against modern FPV systems. Combat footage reviewed for this research indicates that find--engage cycles are measured in seconds\index{Kill-Chain!3 seconds for find--engage}. For planning purposes, sub-minute engagement timelines will outpace human reaction. Hence hesitation or indecision can be fatal. This is reflected in \textcite{Filatov_fpv}, who highlights that improvements in sensors and operator proficiency have improved the detection of stationary targets. Consequently, dispersion, rapid movement and the exploitation of terrain are essential immediate actions for survival \parencite{civdiv_11}.

This creates a fundamental tension in infantry behaviour: movement increases the risk of detection, yet immobility increases the risk of engagement once detected. Commanders must continuously balance exposure against survivability, making rapid decisions on whether to halt or manoeuvre based on incomplete information. The emphasis on SLLS\index{SLLS}\index{Drone Halt} halts reflects an attempt to avoid detection and the dilemmas arising if observed --- producing movement akin to `island hopping'.

Finally, the variation in drone types further complicates this tactical problem. The persistence and altitude of fixed-wing ISR systems mean that infantry must often assume continuous observation. The shorter range of rotary systems creates opportunities to exploit their poor endurance and even to locate and target drone operators \parencite[6:52--13:45]{garand_thumb_ti_2025_2_27}. Indeed, waiter drones have been used to triangulate FPVD teams by watching other drones' flight paths. This suggests that drone avoidance is not purely defensive: in certain conditions, it enables a reversion to traditional infantry strengths, including stalking and counter-reconnaissance. Survival is therefore not about avoiding drones altogether. It is about reducing detectability long enough to avoid being pulled into the kill-chain \parencite[29]{neron_bancel_2024_05_27}.


Drones make routine movement along previously proven routes increasingly hazardous\index{Route Planning}\index{Movement!Route planning}. For example, waiters\index{Waiters} may ambush such routes, while CASDs may drop landmines and IEDs. Hence, the concept of long-term ``proven routes'' becomes increasingly obsolete\index{Proven Route}. Ukrainian SOF practice has institutionalised this logic through structured pre-movement planning, integrating aerial overwatch, dedicated CUAS groupings and continuous route reassessment as baseline requirements rather than optional enhancements \parencite{uKR_LESSONS_LEARNED_1,uKR_LESSONS_LEARNED_2}. Route selection and adaptation cannot reside with higher command; the tempo, decentralised nature and unpredictability of drone-enabled interdiction push these decisions downward to the lowest level. Given individual ownership of risk during movement, Ukrainian units performing infiltration are reported to adopt a decentralised, collaborative planning approach in which the full platoon participates in the estimate rather than receiving orders from above \parencite[11:47--15:00]{Leach_curtain} --- a dynamic examined in Chapter 2.





Under pervasive ISR, the protective qualities of speed and dispersion are emphasised \parencite[62--63]{ZABRODSKYI_2022}. \textcite[65]{mingus_monility,neron_bancel_2024_05_27} reduce this logic to a blunt rule: ``mobility equals survivability'' in a manner at odds with recent Western approaches to protection. The practical expression of this is visible in the proliferation of light mobility platforms --- such as motorcycles and all-terrain vehicles --- appearing on the battlefield. Yet speed alone does not confer immunity. For example, \posscite{Filatov_bike} footage captures the precarious nature of grey zone movement, narrowly avoiding an FPVD strike even when moving at speed. Mobility raises the cost and difficulty of targeting but the significant risk to infantry remains. It creates tensions with the logic of being uninteresting, as the speed and signature required for movement may themselves attract attention. The margin between survival and attrition is often measured in reaction time rather than distance --- especially since vehicles cannot outrun FPVDs. Hence, movement must be aggressive and decisive. However, speed introduces its own hazards: high-tempo movement across broken terrain frequently results in crashes, loss of control and personnel being thrown from vehicles. Speed is therefore not a safeguard, but a trade-off, reducing exposure to the enemy kill-chain.


Thermal sensing further complicates drone avoidance by challenging assumptions about concealment. In terrain permissive to ISR, concealment at night can be more difficult than during daylight, as reduced ambient heating and background clutter increase the detectability of human heat signatures \parencite[2:15--3:00]{civdiv_11}. This inverts the traditional assumption that night movement is inherently safer, with thermal sensors in some cases rendering nocturnal manoeuvre more vulnerable than daylight movement. Even attempts at thermal camouflage may prove counterproductive, as coverings can generate distinct silhouettes or re-radiate absorbed heat following contact. Under such conditions, daylight movement may present lower detection risk, particularly where visual concealment remains effective and detection relies more heavily on movement than static signature \parencite[8:14--8:20]{civdiv_6}. While low-cost mitigation measures such as scrim netting and improvised overhead cover do not eliminate thermal vulnerability, they can introduce friction into the detection process, delaying the transition from observation to engagement \parencite[0:16--50]{civdiv_11}.




%\section{Functional Role Differentiation: Defenders and Assaulters (66 words)}
A functional divergence has emerged between endurance-oriented defenders in dispersed hardened positions and highly-trained assault elements conducting small-team infiltration. This reflects adaptation rather than doctrinal preference. \textcite[8]{WATLING_2023A} note a wide variety of ability amongst Russian infantry; poorly trained soldiers are used for infiltration and reconnaissance by fire while experienced and more capable troops are preserved for assaulting.

%\section{Conclusion (130 words)}

UxS have fundamentally altered the conditions under which light infantry fight, survive and move. The battlespace is no longer opaque; movement has become \textit{timed exposure}; the grey zone has replaced the traditional close--deep--rear model with overlapping zones of opportunity, contestation and risk. UGS extend this logic into sustainment, displacing human movement from the most lethal ground and enabling persistence where exposure would otherwise be prohibitive. Yet these adaptations bring their own constraints: dispersion preserves short-term survival but erodes mass and tempo. Signature discipline and route planning become constant tactical disciplines. The result is a trade-off between survivability and combat power. Light infantry remain indispensable for clearing and consolidating ground, but they now function primarily as survivable shaping and fixing elements within a wider human--machine system. The next chapter examines how these conditions reshape mission command and force design at the lowest tactical level.



%UNUSED IDEAS


%%\section{Tunnels}

%\textcite{civdiv_2} describes his experience with the \gls{ypg} sub-subterranean tunnel complexes which were $5\;--\;500m$ below the surface, $10\,s$ of kilometers long, with multiple exits, rooms and blast doors. RC UAS are ineffective to penetrate such tunnels, due to signal attenuation --- limiting their (and indeed bombers) use to entrances or known bunkers. Common  UAS frames such as $10''\;\ge$, are often too large to traverse tunnels effectively. While he assesses that only AI-piloted UAS may effectively traverse tunnels, they often require GPS connection. While ground drones offer considerations, they utility is also limited.




%%\section{Traumatic Brain Injury}
%Whilst formerly somewhat taboo discussions regarding \gls{tbi} have become commonplace in sport. The military has been slow to follow suit --- arguably concerned with litigation. The prevalence of FPV-OWADs\index{TBI} introduces an additional consideration regarding the long-term effects of combat. Exposure to overpressure of allied force during GWOT --- particularly those exposed to IEDs but also artillery crewmembers, explosive breachers and even SOF with substantial small-arms range-time --- has attracted discussion. That UAS are prevalent, increases the exposure to small high-explosive charges --- under 1kg \parencite[2:58--4:11]{civdiv_3}. While perhaps moot during existential peer-on-peer conflict, the long term medical impact on professional soldiers exposed to FPVDs\index{First Person View Drone!Medical impact of TBI} should not be discounted.




%\textcite[4:06--5:40]{lyle_23_interview} ``when you get into an urban fight, the nature of hte terrain breaks up battles. you odn't have the same capacity to do brigade level manicures or battle group level activities. the largest level unit manoeuvre element you'll see resourced is company group and below. ''tends to break battles down into a much more  small unit fight. you then also end up some of the challenges  that imposes on junior leaders in particular.'' section and platoon commanders have to be ``sources of information for higher formations in a way they're not necessarily required to be in a more spread out rural battle. there's a much greater level of responsibility put on those junior leaders. it's a more complicated terrain they have to contend with. they're in much closer proimity to hte enemy so there's an increased threat to them specifically. they're now also having to coordinate both the close fight and feed information back up. [...] much larger cognitive load, cognitive burden placed on junior leaders at that lowest tactical level 

%\textcite[19:30--21:50]{lyle_23_interview} argued that new technologies are often pushed to platoon and section level without the personnel or institutional adaptation required to exploit them. Faced with competing immediate tasks, soldiers tend to prioritise core defensive functions --- covering their arcs rather than operating drones --- with the result that such technologies are not readily habitualised into tactical practice.

%\textcite[21:50--25:00]{lyle_23_interview} outlined a ``next generation'' combat team in which a company would consist of two larger phalanx manoeuvre platoons\footnote{Comparable logic can be seen in the US Army Ranger platoon, which combines three manoeuvre elements with a support element.} and a manoeuvre support group, the latter permanently integrating capabilities currently held above company level, including GPMG SF, mortars, and Javelin. He emphasised that the key conceptual shift lies in the permanent rather than temporary allocation of these assets to the company. He also noted a tension in relation to combat engineering, arguing that permanently embedded engineer elements would likely take decades to realise and that assault pioneers would therefore have to bridge the gap in the interim\footnote{Such considerations are not trivial. In peacetime, human-resource constraints --- especially promotion pathways and unit cohesion --- would likely stymie the permanent detachment of brigade-level enablers.}.

%drones have a smaller \textit{irreducible minimum} than mortars. Sections can function below full-strength, but the mortar's floor is greater than most. --0 contrast with darraghs' observations that mortarts are the business


%Consider the \textit{[battalion] combat service support group}. just because UGS are robotic, the instinct to `take more risk' mightn't realise because they'll be qiute expensive. \textcite[40:00--41:52]{lyle_23_interview}. recognition that it's not about just assault. the casevac can be more effective with robotics. casevac ugs/uas. `` last mile resupply'' UGS heavy drones. Corroborated by \textcite[18:50--28:00]{taylor_ukr} --- described the SNAKE resupply vehicle with a range of 20km. 


%\footnote{They also highlighted the impact of rubbelisation during peer conflict, where close quarters battle consists of modified combat indicators and the stand-off use of RPGs and hand grenades ``until [there are] no more voices'', corrorobated by depictions such as \textcite{civdivCQB}.}


%\textcite{FOX_2023a} similarly contend that Western doctrine has dogmatically elevated manoeuvre as the sole legitimate mode of warfare, dismissing attrition and positional approaches as obsolete --- 21st centrury warfare demonstrates that these remain contextually valid \parencite{FOX_2021,FOX_2025,Halem_2024}. \textcite{FOX_2021} further argues that manoeuvre requires specific conditions --- open terrain, distance from the enemy, freedom of movement --- and that as contact increases and terrain closes, attrition becomes not a strategic choice but an environmental characteristic. King distinguishes between physical manoeuvre and the conceptual manoeuvrist approach, arguing that precision fires have reshaped the geometry of land warfare such that positional attrition now precedes rather than precludes manoeuvre: infiltration opportunities may emerge once the enemy has been sufficiently attritted, and forces can then seize key objectives through bite-and-hold \parencite[12:58--15:15]{KING_2023}. Consequently, the Ukrainian experience should not displace armoured and air warfare from Western doctrine, but rather inform how attrition and manoeuvre interact under precision-saturated conditions.



%%\section{UAS Considerations, Drone Team Composition and Skills (595 words)}
%Team skills combine radio engineering, tactics, analysis of the EW environment and frequencies and dynamic decision making \parencite[4:45--5:00]{hotoke} \parencite{dan_participant_1,dan_participant_6}. \textcite[13:17--13:17]{zandar_2025} pilot assessed a 10 percent chance of survival for dismounted infantry versus an FPVD. ``We're not manufacturing drones. Every drone that comes to us, has to be modified to fit the mission in some way --- change frequencies, different munitions for example''. \parencite[22:54--25:00]{Ingvald_1}. operate near to the FLOT but within the grey zone. difficult to move to their positoin .require protection, bunker etc. \parencite{Zhluktenko_curtain}
